\chapter{A Research Roadmap}
\label{ch19:top}
\label{ch19}

\section{Where We Stand}
\label{ch19:wherewestand}

Part~I of this book developed the toolkit of analytic combinatorics:
generating functions, singularity analysis, transfer matrices, and
Boltzmann samplers~\cite{FlajoletSedgewick2009}.  Part~II established
the linguistic foundations: weighted finite automata, the Hankel
characterisation, PCFGs, entropy rates, and the Chomsky--Schützenberger
theorem.  Part~III introduced the Transformer architecture and surveyed
the theoretical results bearing on its expressive power and
limitations~\cite{WeissGoldbergYahav2018,MerrillSabharwalSmith2022,RizviEtAl2024}.
Part~IV introduced the global partition function $Z(\beta)$ as a mathematically clean reference object, distinguished it from local decoding temperature, and surveyed a conjectural bridge from combinatorial singularity theory to empirical phase-transition evidence~\cite{BanderierKubaWallner2021}.  Part~V
surveyed Zipf distributions, motif statistics, and their analytic
explanations~\cite{Szpankowski2001}, before Chapter~18 catalogued ten
open problems --- Gaps~1a, 1b, 2, 3, 4, 5, 6, 7, 8, 9, 10 --- at the
frontier.  This final chapter answers a practical question: given that
background, what should a researcher actually work on first?
% NOTE: This chapter works well as a pragmatic roadmap. Since several earlier chapters
% intentionally blended secure theorems with still-speculative LLM bridges, it may help
% to keep signaling here which recommendations depend only on classical mathematics and
% which depend on unresolved conjectural links in the current draft.

\section{Ranking the Gaps by Tractability}
\label{ch19:ranking}

Not all ten gaps are equally open.  Some are computation problems
waiting for careful accounting; others require new mathematics that is
not yet in sight.  The following ranking is a judgement call, not a
theorem.

\textbf{Most tractable: Gap~9 (WFA-approximation pipeline with analytic
error bounds).}  The four approximation strategies surveyed in
Chapter~12 --- spectral learning from moment matrices, AAK-based
Hankel truncation, $L^*$-style automaton learning, and tensor-network
contraction --- are each well understood in isolation.  What is missing
is a unified error budget: given that each step introduces its own
approximation error, what does the compound bound look like?  The
ingredients are already on the shelf.  The AAK singular-value bounds
were made explicit for weighted automata by Lacroce, Panangaden, and
Rabusseau~\cite{LacrocePanangadenRabusseau2021}; the $L^*$ convergence
guarantees are classical.  The missing piece is to compose these bounds
and state a single theorem of the form: ``for a rank-$r$ WFA
approximation to a clearly specified target --- for example, the full string distribution over sequences of length at most
$n$, the total-variation error is at most $f(r,n)$.''  This is, in
essence, a careful accounting exercise.  It does not require proving
anything new about LLMs themselves.

\textbf{Second: Gap~5 (Gaussian limit laws via the Hwang quasi-power
theorem).}  Hwang's 1998 meta-theorem~\cite{Hwang1998} gives Gaussian
limits once the relevant probability generating function satisfies a
quasi-power hypothesis.  Verifying that hypothesis for a
concrete LLM surrogate --- say, a small PCFG with a softmax output head
--- is a calculation, not an open problem in the usual sense.  The
difficulty is making the surrogate tractable enough that the PGF can be
computed or tightly bounded.  A PCFG with $k$ nonterminals has an
algebraic GF of degree at most $k$; the softmax perturbs it in a
controlled way.  Working through this for the simplest non-trivial case
would constitute a publishable result and, more importantly, would
establish whether the quasi-power hypothesis is robust to the softmax
perturbation.

\textbf{Third: Gap~1b (the entropy-rate identity $h = \log(1/R)$ for
LLM counting GFs).}  Once Gap~9 is resolved and a WFA approximation
exists for a clearly specified counting target, Gap~1b becomes a
surrogate-verification problem rather than an immediate corollary.  For
the surrogate itself, Perron--Frobenius gives exact exponential growth
rates for rational counting series.  The real question is whether the
counting object one has chosen is the right analogue of Chapter~9's
entropy-typical construction, and how close its radius-based quantity is
to an independently estimated entropy quantity for the original model.

One must also be precise about the distinction between the length
probability generating function (Gap~1a, concerning $\sum_n p_n x^n$
where $p_n = \Pr[\text{length} = n]$) and a counting GF (Gap~1b,
concerning some sequence $\sum_n c_n x^n$ built from a counting or
entropy-typical object).  These are not the same object and the
singularities can differ.  So the right Phase~2 task is to decide
which counting object is being tested, compute its radius on a
surrogate, and compare the resulting $\log(1/R)$ with the relevant
entropy-like measurement rather than assume the identity automatically.

\textbf{Big prize: Gap~10 (phase transitions as confluent
singularities).}  The empirical observations of Nakaishi et
al.~\cite{NakaishiNishikawaHukushima2024}, Arnold et
al.~\cite{ArnoldEtAl2024}, and
Mikhaylovskiy~\cite{Mikhaylovskiy2025} all point to sharp transitions
in LLM-generated text statistics at critical temperatures.  The
Banderier--Kuba--Wallner composition-scheme framework~\cite{BanderierKubaWallner2021}
gives a rigorous mechanism for such transitions: they arise when two
dominant singularities of the partition function $Z(\beta, x)$ coalesce
as $\beta$ varies.  A positive result connecting these two bodies of
work --- proving that the observed transitions correspond to singularity
confluence in $Z(\beta, x)$ for an explicit surrogate model --- would
be a significant contribution.  This is the right long-term target
but not the right starting point.
This phase also inherits the Chapter~14--15 distinction between local
decoding temperature and the global Gibbs parameter of $Z(\beta)$; any
successful bridge must make that relationship explicit rather than
silently identifying the two.

\textbf{Long shot in the sense of ease of entry, not in probability of
success: Gap~4 (Chomsky--Schützenberger test).}  This is an
empirical project.  The idea is simple: sample strings from a
medium-sized open-weight LLM at fixed temperature, fit a counting GF to
the empirical length distribution, and ask whether the dominant
singularity is algebraic of the type $\sim C \cdot n^{-3/2} r^{-n}$
(the generic square-root signature of a non-degenerate context-free-like
system).  Anyone with
a GPU cluster and access to an open-weight model can start this
tomorrow.  The theoretical interpretation of a positive or negative
result is less straightforward, but the empirical finding would itself
be informative.  This gap is listed last in tractability not because it
is hard to attempt but because the result is harder to interpret without
the theoretical scaffolding of Phases~1--2.
An empirical $n^{-3/2}$ fit here would be suggestive, not a direct proof
that the support language is approximately context-free.

\section{A Three-Phase Programme}
\label{ch19:threephase}

The following programme is designed so that each phase produces
publishable intermediate results while also building the substrate for
the next phase.

\paragraph{Phase~1: Implement and benchmark a WFA-extraction pipeline.}
\label{ch19:phase1}

Choose a small LLM surrogate: a 2-layer Transformer trained on a corpus
generated by a known PCFG.  Using the PCFG, the ground-truth
distribution is computable, so errors can be measured exactly.  Fix the
primary approximation target up front: the most natural first choice is
the truncated full string distribution on strings of length at most $L$.
Extract
the Hankel matrix $H$ of the surrogate's string probabilities up to
length $L$, compute its singular-value decomposition, and truncate to
rank~$r$.  Record the singular-value decay as a function of $r$ (this
is a diagnostic of how well a WFA can approximate the model).  Construct
the rank-$r$ WFA by the spectral-learning procedure of
Balle--Mohri or equivalently via Lacroce--Panangaden--Rabusseau~\cite{LacrocePanangadenRabusseau2021}.
Measure the empirical KL divergence between the WFA distribution and the
surrogate distribution as a function of $r$ and $L$.  Compare to the
theoretical AAK bound.

Deliverables: a codebase, a table of singular-value plots, and a
graph of empirical KL divergence versus rank.  The gap between the
empirical error and the AAK bound, if large, is itself an interesting
finding.

\paragraph{Phase~2: Verify the entropy-rate identity.}
\label{ch19:phase2}

Using the rank-$r$ WFA from Phase~1, compute the transition matrix $M$
and extract $\rho(M)$.  Compute the exact rational objects associated
with the surrogate's chosen counting target, including the relevant
radius $R$, and compare $\log(1/R)$ to an independently measured
entropy-like quantity on the surrogate (for example, a cross-entropy
estimate on held-out strings).  Do this
for several values of $r$ and record how quickly the discrepancy
decays.

This does not automatically close Gap~1b for the original LLM. What it
does provide is a controlled testbed: which counting object supports a
log-radius identity on the surrogate, and how sensitive that identity is
to approximation error.

\paragraph{Phase~3: Scan temperature and look for singularity
confluence.}
\label{ch19:phase3}

Fix the surrogate from Phases~1--2.  Define the temperature-parameterised
finite-size proxy
\[
Z_L(\beta) = \sum_{|\mathbf{x}| \le L} p(\mathbf{x})^\beta,
\]
where the sum runs only over strings up to length $L$.  For each
$\beta$ in a grid, compute the dominant singularity of the associated
surrogate WFA (the reciprocal spectral radius $R(\beta)$) and track how
it moves with $\beta$.  This is a finite-size probe, not yet a literal
nonanalyticity of the infinite-volume partition function.  A sharp kink
or crossover in $R(\beta)$ is evidence for a singularity crossing, which
in the
Banderier--Kuba--Wallner framework~\cite{BanderierKubaWallner2021}
corresponds to a phase transition.  If a transition is found, attempt to
classify the type of confluence (smooth, cusp, or higher order) by
examining the local expansion of the surrogate partition object near the critical
point.
In parallel, one may scan the empirically implemented local decoding
temperature on the teacher model and compare observables, but the two
temperature operations should not be conflated.

A positive finding here constitutes empirical evidence for Gap~10.  A
negative finding is also informative: it would suggest that LLM phase
transitions, when they occur, are not of the composition-scheme type
and would require a different theoretical explanation.

\section{Parallel Tracks and What to Avoid First}
\label{ch19:parallel}

Gap~4 (the Chomsky--Schützenberger numerical test) is an orthogonal
project that does not depend on the WFA pipeline.  It can be pursued
in parallel by a researcher with access to a real open-weight
LLM~\cite{BlondelEtAl2025,DuEtAl2023} and a good sampling implementation.
The main risk is that the length distribution of LLM output is heavily
confounded by prompt length and stopping heuristics; controlling for
these requires care. Another confounder is the local-vs-global
temperature distinction from Chapter~14: fitting a global analytic model
to tails produced under local decoding temperature can be conceptually
ambiguous.

Gap~5 (Hwang quasi-power limits) can be pursued as a theoretical
side-project by a combinatorics-oriented researcher who does not have
access to GPU resources.  The calculation for the PCFG-plus-softmax
surrogate can be carried out on paper or with a computer-algebra
system~\cite{Hwang1998}.

\textbf{What not to work on first.}  Three gaps should be deferred.

Gap~6 (thermodynamic formalism) is where the cleanest theory lives:
the Ruelle--Bowen--Sinai transfer operator gives a complete account of
the free energy and its singularities.  But this background is a
semester course on its own, with prerequisites in ergodic theory that
most combinatorics or machine-learning students will not have.  Return
to it after Phases~1--3 have built intuition.

Gap~7 (classification of all singularity types arising in LLM partition
functions) is probably the ultimate theoretical destination of this
research programme, but it presupposes Gap~10 being resolved and a
rich collection of examples to classify.  It is not a starting point.
It also presupposes deciding what analytic class the partition object
$Z(\beta)$ actually belongs to.

Gap~3 (ACSV for tokenized models) requires the Pemantle--Wilson
multivariate machinery~\cite{PemantleWilson2013} in its full
generality.  This is technically demanding --- the critical-point
analysis for multivariate GFs involves stratified Morse theory ---
and the payoff depends on first establishing that LLM generating
functions have the right structure to make the analysis applicable.
Defer until the univariate picture (Phases~1--2) is settled.

\section{For a Student Entering the Field}
\label{ch19:student}

The following reading order is recommended for a graduate student
who has not yet read this book and wants to enter this research area
efficiently.

\begin{enumerate}
  \item \textbf{Flajolet--Sedgewick~\cite{FlajoletSedgewick2009},
  Chapters~I--IV.}  This covers symbolic combinatorics, ordinary and
  exponential GFs, and the singularity-analysis framework.  These
  chapters correspond to Chapters~1--5 of this book.  This material is
  non-negotiable; everything else builds on it.

  \item \textbf{This book, Part~II} (WFAs, PCFGs, entropy rates).
  The primary mathematical objects are the Hankel matrix and the
  Perron--Frobenius theorem applied to weighted transition matrices.

  \item \textbf{Szpankowski~\cite{Szpankowski2001}.}  For researchers
  whose interest is in word statistics, motif occurrences, and
  large-deviation results.  Chapter~17 of this book is an appetiser;
  Szpankowski's monograph is the full treatment.

  \item \textbf{Rizvi et al.~\cite{RizviEtAl2024} and
  Merrill--Sabharwal--Smith~\cite{MerrillSabharwalSmith2022}.}  For the
  LLM side: what Transformers can and cannot express as formal languages,
  and what circuit-complexity arguments say about their computational
  depth.  Read alongside Strobl et al.~\cite{StroblEtAl2024} for the
  broader survey of formal language results.

  \item \textbf{Duchon--Flajolet--Louchard--Schaeffer~\cite{DuchonFlajoletLouchardSchaeffer2004}.}
  The foundational reference for Boltzmann sampling.  Read this before
  attempting Phase~3; the connection between the Boltzmann parameter and
  the thermodynamic $\beta$ is made explicit there.
\end{enumerate}

This reading list is not comprehensive, but it covers the minimum
required to engage with the three-phase programme above without
repeatedly needing to pause and build background.
Chapters~1--8 and the exact finite-state parts of Chapters~12 and~17 are
best read as classical or exact toolkit. Chapters~9, 14, and 15 are more
bridge-like and conjectural in places; they are important, but they
should be read with that status in mind.

\section{Closing}
\label{ch19:closing}

The purpose of this book has been to place in one document the pieces
of two literatures --- analytic combinatorics and large language model
theory --- that have developed largely independently.  The mathematical
community has Flajolet--Sedgewick, Szpankowski, and Pemantle--Wilson;
the machine-learning community has Transformer theory, scaling laws,
and empirical phase-transition observations.  These are not separate
subjects.  The formal-language perspective forces precision about what
a language model is computing; the singularity-analysis perspective
gives a vocabulary for what happens at critical temperatures.  If the
book has done its job, a reader should be able to move between a
chapter of Flajolet--Sedgewick and a survey such as Strobl et
al.~\cite{StroblEtAl2024} without feeling that they have changed
disciplines.  The ten gaps of Chapter~18 and the programme sketched
here are an invitation to close that distance further.  They are not
a manifesto; they are a list of things that seem doable.

\section*{Further Reading}
\label{ch19:furtherreading}

The reading list in Section~\ref{ch19:student} covers the minimum
required for the three-phase programme.  The items below are the core
references specifically for the roadmap itself.

\begin{itemize}
  \item \textbf{P.~Flajolet and R.~Sedgewick}, \emph{Analytic Combinatorics} (Cambridge University Press, 2009) \cite{FlajoletSedgewick2009} --- The analytical backbone of Phases~1--3. Chapters~VI--IX on singularity analysis, algebraic functions, and the saddle-point method are used directly; Chapter~IX on words and tries is the template for the WFA transition-matrix computations in Phase~2.

  \item \textbf{W.~Szpankowski}, \emph{Average Case Analysis of Algorithms on Sequences} (Wiley, 2001) \cite{Szpankowski2001} --- The bridge between analytic combinatorics and information theory. For researchers entering from the combinatorics side, this monograph makes entropy rates, trie depths, and suffix-tree statistics into singularity-analysis problems, exactly the translation needed for Phase~2.

  \item \textbf{R.~Pemantle and M.~C. Wilson}, \emph{Analytic Combinatorics in Several Variables} (Cambridge University Press, 2013) \cite{PemantleWilson2013} --- Required for Gap~3 and for Phase~3 once the analysis extends beyond univariate generating functions. The critical-point and stratified-Morse-theory toolkit for multivariate coefficient asymptotics.

  \item \textbf{H.-K. Hwang}, ``On convergence rates in the central limit theorems for combinatorial structures,'' \emph{European J.\ Combin.} (1998) \cite{Hwang1998} --- The quasi-power meta-theorem that converts a bivariate generating function in quasi-power form into a Gaussian limit law. A short paper with large returns: reading it clarifies exactly what Phase~1's WFA extraction must deliver for Gap~5 to follow automatically.

  \item \textbf{D.~Ruelle}, \emph{Thermodynamic Formalism} (Addison-Wesley, 1978) \cite{Ruelle1978} --- Background for Gap~6 and for understanding why Phase~3's $Z(\beta)$ behaves like a thermodynamic pressure function. The Ruelle transfer operator is the infinite-alphabet generalisation of the WFA transition matrix.

  \item \textbf{R.~Bowen}, \emph{Equilibrium States and the Ergodic Theory of Anosov Diffeomorphisms} (Springer, 1975) \cite{Bowen1975} --- Companion to Ruelle for the variational principle and Gibbs-state characterisation. Together these two short books give the ergodic-theoretic language that the roadmap's Phase~3 is, in spirit, trying to instantiate for neural-network potentials.
\end{itemize}
